%&../template
\documentclass[../main.tex]{subfiles}

\begin{document}

    \clearpage
    \section{Introduction}

    In this Thesis, we want to explore the standard kernel User I/O (\UIO) framework and discuss possibilities of integrating the FIQ Interrupt on ARM \blindCite. We will be using the \href{https://www.raspberrypi.com/products/raspberry-pi-4-model-b/specifications/}{Raspberry Pi 4 Model B Rev 1.5} as it contains all vital pieces of hardware to do the benchmarking correctly and accurately.

    In order to cite the use of AI properly, it is integrated as a \enquote{Live}-Chat into the thesis. Whenever it was relevant to ask the AI, the \textcolor{Orchid}{Question} and \textcolor{red}{Answer} (word-for-word or excerpt) will be provided and shown. Blackbox.ai was used for providing the \textcolor{red}{GPT-5.5} Model.

    \AICite{Hey, I'm new to this. I don't know what the FIQ Interrupt is on ARM and what Benefits it has in comparison towards a \enquote{normal} IRQ in x86-64 / arm64?}{FIQ = Fast Interrupt reQuest. It is an ARM interrupt type intended for very low-latency, high-priority interrupt handling. \AIParen{} FIQ is not really an x86-64 concept, it has normal maskable interrupts, NMIs, SMIs, APIC priority levels, etc., but there is no exact direct equivalent to ARM FIQ.}
    %My Answer: {The FIQ Interrupt is of particular note as it is \enquote{reserved for a single, high-priority interrupt source that requires a guaranteed fast response time}\footnote{https://developer.arm.com/documentation/den0013/d/Exception-Handling/Exception-priorities/FIQ-and-IRQ}.}
    %

    \textsc{Now}, the interesting question becomes, \textbf{how} does ARM do this?\lf
    \textit{By implementing the entire Architecture from scratch and introducing this concept.}\lf
    This is only possible because the Architecture (\texttt{ARM-7} / \texttt{RISC-V}) isn't clogged by decades of legacy.


    \bigskip


    \subsection{\GPT5 and the \ARM{} Architecture}

    We train our specialized \AI{} model with \GPT5 and the following data

    \begin{itemize}
        \item \CPU: \Cortex
        \item \MEM: 8 \GB
        \item \MAC: \verb|d8:3a:dd:46:f2:81|
    \end{itemize}

    \url{https://developer.arm.com/documentation/ddi0487/gb/}

    \url{https://pip-assets.raspberrypi.com/categories/545-raspberry-pi-4-model-b/documents/RP-008248-DS-1-bcm2711-peripherals.pdf}

    All of the files in \url{https://github.com/Emily3403/Bachelor\_Thesis/tree/main/resources} are relevant for the Large Language Model of the AI. For technical specificies, only the \verb|broadcom| directory was used. For everything else, the AI – STEVE <bauroth@tu-berlin.de> –  was trained on the data around this directory.

    \TODO[Other sources]


    ---

    Trained with this data, \GPT5 is able to answer all questions related to the Raspberry Pi 4 B and the \ARM{} architecture. Let's call them Stephan [Bauroth <\url{s.bauroth@tu-berlin.de}>].

    
    \subsection{New}



    % Arguments:
    %   - Higher Prio, won't be disabled
    %   - Banked registers
    %   - IVT entry at end → no branch
    %   - Only one interrupt source, no need to figure out which device caused the interrupt

    This property is achieved in two steps: Firstly, the FIQ interrupt has a higher priority than IRQ and won't be disabled when handling other exceptions: \enquote{All exceptions disable IRQ, only FIQ and reset disable FIQ}$^1$. Secondly, the ARM Architecture provides an additional 5 banked registers specific to the FIQ mode, allowing us to retain 40 bytes of memory across interrupts. Moreover, the FIQ entry in the Interrupt Vector Table (IVT) is the last one, making it possible to place the handler code directly afterward and skipping a branch instruction and the associated delay. With all these \todo[tricks], we are able to shave off precious processing time and decrease latency. In Chapter \todo, we will be measuring the concrete performance difference between FIQ and IRQ with the same driver code, just with swapped out interrupt implementations.

    One important caveat of the FIQ is, that the handler is not able to generate executions such as a supervisor calls (SVC). This however integrates very nicely with the UIO framework, as \enquote{only a small kernel module is needed for handling the interrupt}, reducing the likelihood of needing any SVCs. Once the interrupt arrives in userspace, the ability to call SVCs is restored, so we won't be restricted by this limitation. Further, the FIQ interrupt is not typically used by Linux making it ideal to map the single interrupt source into userspace and increase performance for one specific userspace application.


    \clearpage

    With more and more drivers in the Linux Kernel switching to userland, we want to investigate how to improve performance with these drivers.

    Writing drivers in userland has numerous advantages such as
    \begin{itemize}
        \item Wrong code won't crash the entire system
        \item No kernel invariants can be messed with
        \item Faster Debugging
        \item Easier Devolopment
    \end{itemize}

    On the contrary however, these drivers also suffer from a few caveats introduced by not having the code run in kernel space:
    \begin{itemize}
        \item Performance overhead due to context switches
        \item Memory Limitations
        \item Userspace threads might be swapped out
        \item During memory pressure, the kernel might decide to kill the userspace driver
        \item Application restart: If the application is not given the opportunity to clean up or crashes, the device might be left in an inconsistent state
    \end{itemize}

    % Further Content:
    %   - Warum Raspi 3b?
    %
    %   - IRQ Paper
    % - Userspace can't get hardware interrupts directly → no routing directly to user-level applications. Maybe this is an assumption I could change with my kernel driver?
    % - Userspace isn't integrated with scheduler
    % - IRQs distort the processor pipeline quite significantly
    % - Why doesn't DPDK use Interrupts instead of polling? Why is the overhead so high?
    %
    % Practical considerations
    % - Interrupt Routing: which CPU core(s) handle my interrupts?

    In this Thesis, we will primarily focus on the Mini-UART (\texttt{UART1}) found on the bcm2711 chip. This is because of two reasons: Firstly, it contains shallow FIFOs (8 bytes) that have to be read by the CPU and can't benefit from hardware acceleration or DMA. Secondly, it can run up to 32 MegaBaud.


\end{document}
